\documentclass[12pt]{article}
\usepackage{amsmath, amssymb, amsfonts, geometry, setspace, tcolorbox, amsthm}
\geometry{margin=0.75in}
\setstretch{1.3}
\setlength{\parindent}{0pt}

\tcbset{colback=gray!5, colframe=black!70, boxrule=0.6pt, arc=2pt, 
left=6pt, right=6pt, top=2pt, bottom=2pt}

\newtheorem{definition}{Definition}
\newtheorem{theorem}{Theorem}
\newtheorem{example}{Example}

\title{\textbf{Introduction to Behavioural Economics}\\[6pt]
\textbf{ECON 230}}
\author{}
\date{}

\begin{document}

\maketitle

\section*{Course Description}

This course explores why individuals, firms, and public institutions sometimes make choices that seem careless, puzzling, or harmful. Drawing on theoretical and empirical perspectives from economics and psychology, the course examines how temptation, bias, social influence, and strategic behaviour lead decisions away from the predictions of standard rational models. Topics include financial fraud, addictive consumption, tax evasion, workplace misconduct, and deceptive business practices,  analysed through real-world cases and applied examples.

Students will learn how habits, short-term thinking, and mistaken beliefs shape behaviour, and how these tendencies can spread and persist in markets and organizations. By the end of the course, students will be able to apply behavioural concepts to explain harmful or irrational choices, assess when markets amplify or limit these behaviours, and evaluate the effectiveness of policy tools aimed at reducing them. The course emphasizes clear reasoning, evidence-based analysis, and critical thinking about how people respond to incentives and psychological forces.



\section*{Course Description (Short Version)}
This course examines why individuals, firms, and institutions sometimes make harmful or irrational choices. Using theoretical and empirical perspectives from economics and psychology, students explore how biases, habits, and social influences shape behaviour in areas such as financial fraud, addiction, tax evasion, workplace misconduct, and deceptive business practices. The course also analyzes the conditions under which markets amplify or mitigate these behaviours and evaluates the effectiveness of policy interventions designed to address them.  Emphasis is placed on evidence-based analysis and a systematic assessment of the limits of standard rational-choice models.

\end{document}
